% Optional: use minted for code listings
%\RequirePackage[]{minted}

\documentclass[
  a4paper,        % select between `a4paper' and `letterpaper'
  techreport,     % select between  `techreport', `report', `article', `commun', `persp' and `review'
  % lineno,         % switch, adds line numbering
  % easyreview,     % use easyReview package
  % anonymous,
  % noauthors,
  % nodate,
]{art}


% Title
\newcommand{\pubtitle}{Mailbox Actors}

\newcommand{\pubauthA}{Jonathan Prieto-Cubides}
\newcommand{\pubaffilA}{a}
\newcommand{\orcidA}{0000-0002-8449-3812}
\newcommand{\authemailA}{jonathan@heliax.dev}
% \newcommand{\eqcontribA}{}

\newcommand{\pubauthB}{Anthony Hart}
\newcommand{\pubaffilB}{a}
\newcommand{\orcidB}{0009-0004-4116-5538}
\newcommand{\authemailB}{anthony@heliax.dev}
\newcommand{\eqcontribB}{}

\newcommand{\pubauthC}{Tobias Heindel}
\newcommand{\pubaffilC}{a}
\newcommand{\orcidC}{0000-0003-3371-8564}
\newcommand{\authemailC}{tobias@heliax.com}

% Institutions/Affiliations
\newcommand{\pubaddrA}{Heliax AG}

% Corresponding author mail
\newcommand{\pubemail}{\authemailA}

\newcommand{\pubabstract}{%
%
We present a formal framework for actor systems in which mailboxes are promoted
to independent, first-class actors. In contrast to traditional actor models,
where mailboxes are treated as opaque runtime primitives or passive data
structures, our design decouples message storage from processing logic by
assigning each processing actor a dedicated mailbox actor. This separation of
concerns enables flexible message scheduling, modular mailbox implementations,
and independent scaling of message ingestion and processing.

We formalise the system using a predicative type theory with dependent types,
giving precise definitions of engines---our refinement of the actor
notion---together with their configurations, environments, behaviours, and
effect systems. We provide a complete operational semantics expressed as a
labelled transition system, covering engine creation, message passing, behaviour
evaluation, effect execution, and engine lifecycle management. We establish
several metatheoretic properties of the system, including type preservation
under reduction, progress, effect determinism under non-overlapping guards, and
mailbox isolation.
}


% Any keywords to be displayed under the abstract
\keywords{
Actor Model \sep
Message-Passing Systems \sep
Mailbox Actors \sep
Operational Semantics \sep
Dependent Types \sep
Formal Methods
}

% Supplementary space between title/abstract and text, if needed
% \newcommand{\pubVadj}{0pt}

% ! DO NOT REMOVE OR MODIFY !
\input{aux-preamble.tex}
% The preprint DOI to be used as a link in the paper
\pubdoi{10.5281/zenodo.14915469}
% \history{(Received Month DD, YYYY; Revised Month DD, YYYY ; Published: August 21, 2023)}
\pubdate{\today}

% MACROS
\newcommand{\N}{\mathbb{N}}
\newcommand{\execEffect}{\ensuremath{\mathsf{execEffect}}}
\newcommand{\engine}{\ensuremath{\mathsf{engine}}}
\newcommand{\engineBehavior}{\ensuremath{\mathsf{engineBehavior}}}
\newcommand{\engineConfig}{\ensuremath{\mathsf{engineConfig}}}
\newcommand{\engineEnv}{\ensuremath{\mathsf{engineEnv}}}
%msg
\newcommand{\Msg}{\ensuremath{\mathsf{Msg}}}

%\input{templates/pandoc.tex} % needed for pandoc-generated files

\begin{document}
\maketitle

\tableofcontents
\section{Introduction}

% Introduction to actor models and their communication mechanism.
Actor models provide a computational framework for concurrent systems. In these
systems, the main construct is the notion of an \emph{actor}. Actors are
independent computational units that encapsulate the execution environment of
their processing logic. Asynchronous communication between actors occurs
solely via messages. A message is defined as any piece of data sent from one
actor to another. The reception of these messages is handled by the recipient
actor's \emph{mailbox}, a dedicated piece of storage that preserves message
order while remaining separate from the actor's core logic.

% Challenges in actor system message flow management
One challenge in actor systems is managing message flow under high volumes.
Although some techniques such as rate limiting and \emph{backpressure} help regulate
message rates, these control mechanisms may also fail under heavily loaded
systems. Message delivery is often assumed to be reliable, but practical
implementations really need to consider the possibility of message loss and
delays from processing bottlenecks or crashing actors.

% Proposes a solution by decoupling mailboxes from processing actors.
To address some of these challenges, in particular the need to handle message
loss and delays in processing messages, we propose decoupling mailboxes
from \emph{processing actors}, which are the actors that carry out the processing
logic. In this design, communication occurs through message passing to
dedicated \emph{mailbox actors}, with processing actors pulling messages only
when ready to handle them. This approach is simple, allowing us to give an explicit
specification of the message-passing mechanism, where the mailbox actors remain
opaque to both senders and receivers to maintain compatibility with existing actor
frameworks.

% Outlines the benefits of the proposed architectural separation.
The architectural separation of mailboxes from processing actors provides
several key benefits. First, it cleanly separates message management from
processing logic, allowing mailbox actors to specialize in queue management
while processing actors focus on their core functionality. Second, it enables
flexible message scheduling—different mailbox implementations can support
priority queues, batch processing, or custom dispatch strategies without affecting
processing actors. Third, mailbox actors become modular components that can be
modified or replaced independently.

% Discusses the trade-offs and potential drawbacks of the design.
This design may come with some trade-offs. The primary drawback we foresee is
increased system complexity due to additional actors and message-passing
overhead. This manifests as higher resource usage from running more actors and
potential latency increases when mailbox actors run on different nodes than
their processing actors. However, with careful design, these costs can be
outweighed by the benefits of a more flexible and scalable system.

% Mentions the potential for future work in extending the design.
The design could be extended to support multiple mailbox actors per processing
actor, with minimal additional complexity due to the mailbox actors' simplicity
and well-defined message interface. However, we leave this extension as future
work since it is not required to convey the core design.

% How this paper is organized and what is the contribution.

\paragraph{Contributions}
First, we propose a new architectural
separation between processing actors and their mailboxes in \Cref{sec:model},
which enables improved flexibility and scalability for message management. Second,
we formulate operational semantics in \Cref{sec:semantics} for both the decoupled
architecture and its explicit message-passing mechanism, making the system's
communication patterns formally explicit. This description provides sufficient
grounding to implement a library on top of a compatible system/platform such
as Erlang, and to formalize and prove key system properties such as liveness
guarantees and correctness under concurrent message processing.

\section{Notation}\label{sec:notation}

We use an informal type-theoretic notation. For example, a term $x$ of type $A$
is denoted as $x : A$. Capital letters or capitalized words are used for types,
and lowercase letters or lowercase words are used for terms. There is a universe
of types $U$, so if $A$ is a type, then $A : U$. We also work with inductive
types presented in the following form:

\begin{align*}
  A ::= \cdots \mid \mathsf{C}(A_1, \dots, A_n) \, ,
\end{align*}

where $\mathsf{C}$ is a constructor of the type $A$, and $A_1, \dots, A_n$ are
the types of the constructor arguments. The inductive type of natural numbers is
denoted as $\N$. The terms of this type are denoted by number as usual. Another
recurrent type is the type of parametric lists, denoted as $\mathsf{List}(A)$.
The definition that term $x$ \emph{is} $y$ is denoted as $x :\equiv y$.
Definitionally equal terms are denoted as $a \equiv b$.

Some types are equipped with structure that behave as sets, and here are common
set-theoretic notation we use: the membership relation is denoted as $\in$, so
$a \in A$ is a proposition that $a$ is an element of $A$ (do not confuse with
typing $a : A$). For $A$ and $B$: $A \subseteq B$ denotes set inclusion, meaning
every element of $A$ belongs to $B$; $\varnothing$ denotes the empty set; union
is denoted as $A \cup B$; intersection is denoted as $A \cap B$; the difference
intersection is denoted as $A \cap B$; the difference between $A$ and $B$ is
denoted as $A \setminus B$; and $A \times B$ denotes the Cartesian product of
$A$ and $B$, and their terms are denoted as $\langle a,b \rangle$, where $a : A$
and $b : B$. A \emph{tuple} $\langle a_1, a_2, \dots, a_n \rangle$ is a
generalisation of an ordered pair to an ordered $n$-tuple, where $n \in \N$, of
type $A_1 \times \cdots \times A_n$, and $a_i$ is of type $A_i$ for all $i \in
\{1, \dots, n\}$. If the tuple components are named using sans-serif letters,
these names are used as projections, otherwise, the projections are $\pi_i(a)
:\equiv a_i$ for all $i \in \{1, \dots, n\}$. For example, if $x :\equiv \langle
\mathsf{fst}, \mathsf{snd} \rangle$ denotes a tuple with components values can
be projected using $\mathsf{fst}$ and $\mathsf{snd}$ via $\mathsf{fst}(x)$ and
$\mathsf{snd}(x)$, respectively. Additionally, $A_1 + \cdots + A_n$ denotes the
disjoint union of $A_1, \dots, A_n$, so $\mathsf{inj_i}(a_i)$ are the
corresponding injections into $A_1 + \cdots + A_n$ from $A_i$ for all $i \in
\{1, \dots, n\}$. Given a type $A$, we can define the optional type $\Option{A}$
as $A + \bot$, where $\bot$ is the bottom type.

\section{Model}\label{sec:model}

This section presents a version based on the Engine Model \cite{engine-art}. In
this actor-like model, the system consists of \emph{engines}, which are similar
to actors but are equipped with more context. An engine comprises three main
components: a \emph{configuration}, an \emph{execution environment}, and a
\emph{behaviour}. The configuration includes, among other things, the engine's
name, location within the system, and an optional reference to the parent
engine. The execution environment includes an \emph{address book} of other engines
to communicate with, including the engine's own address. Finally, the behaviour
defines how the engine reacts to messages that follow the engine's \emph{message
interface}, which is simply a set of message types that the engine can process.
The engine's behaviour is a set of \emph{guarded actions} consisting of pairs
of pre-conditions and actions to execute. When an engine receives a message, it
checks which of its guarded actions have satisfied pre-conditions, and executes
the corresponding action.

\begin{remark}
  The way engine behaviour is defined is similar to the Event-B formalism in
  that event-driven system transitions are specified by a set of pre-conditions
  on sets that form the current system state and the actions to execute if those
  pre-conditions are satisfied. The actions are defined using a language derived
  from process calculi, which often include operations such as parallel composition,
  a null operation, sequencing, nondeterministic choice, and recursion. Here, we
  have a more restricted set of actions, as will be seen in \Cref{def:effects}.
\end{remark}

\subsection{Definitions}\label{sec:definitions}

The model assumes an ambient \emph{system} $\kappa$ consisting of a set of
\emph{nodes} $\mathcal{N}$, a set of \emph{engines} $\mathcal{E}$, following a
specific \emph{message-passing mechanism} (MPM for short) that will be described
later, among other components.

Engines are distributed across the nodes of the system, with each engine being
associated with a unique node. Each entity in the system has a unique identifier,
allowing nodes and engines to be denoted as $n_i$ and $e_j$ respectively, where
$i, j \in \N$.

Following the message-passing paradigm, the MPM enforces policies for message
delivery and flow control. This mechanism operates on a set of \emph{messages}
that we denote here as $\mu$. Each message contains information about its
sender and receiver. 

To reason about the system, we use a state machine that evolves in discrete steps.
This means, for each step/transition of the state machine, we have a new state
$\kappa^{i+1}$ from the previous state $\kappa^i$. This also entails that we
can refer to \say{previous} messages $\mu^i$ or the next ones.

Within the system, each entity has a unique identifier. However, engines can also be
referenced by name rather than by address. Engine names provide a more legible way
to refer to engines in the system, making them easier to work with. Since names
are unique and unambiguous, they can be used interchangeably with addresses when
referring to engines.

\begin{definition}\label{def:address}
An \textbf{address} for an entity $a$ is a tuple $\langle n_i, \mathsf{id}_j \rangle$, where
$n_i$ is the node identifier where the entity lives, and $\mathsf{id}_j$ is the entity's unique identifier.
\end{definition}

\begin{axiom}\label{ax:address}
Each entity is uniquely identified by an \emph{address} in the system.
\end{axiom}

\subsubsection{Engine components}

\begin{definition}\label{def:engine-status}
The \textbf{status} of an engine $e_j$, which is of type
$\mathsf{EngineStatus}$, determines the engine's operational mode and thus its lifecycle.
\begin{equation}\label{eq:engine-status}
\mathsf{EngineStatus} ::= \mathsf{ready} \mid \mathsf{busy} \mid \mathsf{terminated}
\end{equation}
\end{definition}

An engine is said to be \textsf{ready} if it is waiting for a message to
process. Being \textsf{busy} means it is
currently processing a message. Being \textsf{terminated} means it will not
process any more messages.

\begin{definition}\label{def:engine-configuration} The \textbf{configuration} of
an engine $e_j$ often denoted by $\mathsf{cfg}_j$ is a tuple $\langle
\mathsf{name}, \mathsf{location}, \mathsf{parent} \rangle$, where $\mathsf{name}$
is the engine's name, $\mathsf{location}$ is the engine's location, and $\mathsf{parent}$
is the engine's parent.
\end{definition}

Recall engines are isolated from each other, and process messages in isolation.
To properly function, the configuration of an engine is not enough. Engines
encapsulate their own state, as that is part of its execution context, also
called the \emph{environment} of the engine.

\begin{definition}\label{def:engine-environment} 
  Given a type $S_j$ representing the local state of an engine $e_j$, an execution
  \textbf{environment} of $e_j$ is a tuple $\langle s_j, a_j \rangle$, where $s_j$ is 
  the engine's local state of type $S_j$, and $a_j$ is the engine's address book containing 
  name of other engines, including its own. $a_j$ is of type $\mathsf{List}(\mathsf{Name})$, 
  The type of such an environment is denoted 
  as $\mathsf{Env}_j(S_j)$.
\end{definition}


\subsubsection{Execution directives}

\begin{definition}\label{def:effects}
An \textbf{action} is of type $\mathsf{Effect}$, which considers all the
possible ways an engine can modify the system.


\begin{equation}\label{eq:effects}
\begin{aligned} %TODO: fix this grammar
\mathsf{Effect} ::= &\quad \mathsf{send}(\mathsf{Addr}, \mathsf{Msg}) & &\text{(message dispatch)} \\
\mid & \quad \mathsf{update}(\mathsf{Env}) & &\text{(engine state modification)} \\
\mid & \quad \mathsf{spawn}(\mathsf{Cfg}, \mathsf{Env}) & &\text{(engine creation)} \\
\mid & \quad \mathsf{chain}(\mathsf{Effect}, \mathsf{Effect}) & &\text{(action sequencing)} \\
\mid & \quad \mathsf{terminate} & &\text{(Engine shutdown)} \\
\mid & \quad \mathsf{nop} & &\text{(Null operation)\, .}
\end{aligned}
\end{equation}
\end{definition}

% TODO: mention how this induces a monad.


\begin{definition}\label{def:guards} Given an engine $e_j$, a \textbf{guard} is
  a pre-condition that determines how an engine can process a message, including
  ignoring the message, or modifying the system in some way. These guards are
  defined as a mapping from a message, an engine's environment $\mathsf{env}_j$,
  and an engine's configuration $\mathsf{cfg}_j$, and returns an action the system
  must execute if the guard is satisfied.
\end{definition}

Upon receiving a message, an engine $e_j$ checks which of its guards is
satisfied, and if so, the engine executes the corresponding action.




\begin{definition}\label{def:processing-engine}   
A \textbf{processing engine} or \textbf{engine} for short is the sole computational unit of the system, formally specified by a tuple:
\[
\langle b,\, m,\, r,\, \mathsf{cfg},\, \mathsf{env} \rangle
\]
where:
\begin{itemize}
\item $b$ is the engine's behaviour, a set of guarded actions,
\item $m$ is the address of the dedicated mailbox engine,
\item $r$ is the engine's status flag.
\item $\mathsf{cfg}$ is the engine's configuration, a tuple of engine name, network location, and parent reference
\item $\mathsf{env}$ is the engine's execution context, a tuple of local state and address book
\end{itemize}
\end{definition}

\begin{definition}\label{def:mailbox-engine} 
A \emph{mailbox engine} is a specialized processing engine with:
\begin{itemize}
\item Fixed message interface supporting $\mathsf{append}$ and $\mathsf{pull}$ operations
\item Opaque state implementation (queue, priority queue, etc.)
\item Automatic creation with parent processing engine
\end{itemize}

Key responsibilities include:
\begin{itemize}
\item Buffering incoming messages for parent engine
\item Preventing message loss through managed storage
\item Regulating message flow to prevent overload
\end{itemize}

\end{definition}
\begin{remark}
While processing engines have domain-specific behaviours, mailbox engines
maintain fixed semantics for reliable message management. The state
representation remains implementation-defined to support different message
handling strategies.
\end{remark}

\paragraph{Message Structure}

\begin{definition}\label{def:messages}
Engine communication occurs through typed messages where:
\begin{itemize}
\item $\mathsf{Msg}$ denotes the union of all valid message types
\item Messages are triples $\langle s,\, t,\, v \rangle$ with:
  \begin{itemize}
  \item $s$: Sender address
  \item $t$: Target address
  \item $v$: Payload of type $\mathsf{Msg}$
  \end{itemize}
\end{itemize}
\end{definition}

\begin{definition}\label{def:node} 
A \textbf{node} $n_i$ is a tuple $\langle \mathsf{id}_i, \mathcal{E}_i
\rangle$ where $\mathsf{id}_i$ is the node's unique identifier and
$\mathcal{E}_i$ is the set of engines living on the node.
\end{definition}

\begin{remark}\label{rem:node-functions}
Given a node $n_i$, we derive the following functions:
\begin{itemize}
\item $\alpha_i : \mathcal{E}_i \to \mathsf{Addr}$ as the mapping that projects out the addresses of the engines
living on node $n_i$, and
\item $\mathsf{mailboxOf} : \mathcal{E} \to \mathsf{Addr}$ as the mapping that projects out the address of the
dedicated mailbox engine for a given processing engine instance.
\end{itemize}
\end{remark}

\begin{definition}\label{def:global-config}
A \textbf{system state} $\kappa$ is a tuple $\langle \mathcal{N}, \mu, \omega
\rangle$ where $\mathcal{N}$ is the set of nodes, $\mu$ is the in-transit
communications consisting of messages, and $\omega$ is the system capabilities
and operations.
\end{definition}

\noindent It is up-to the discretion of the system designer to define and update
the components of $\omega$. That is, the system can evolve over time, including
the set of nodes, engines, and the included services provided by the system.

We expect at least the following functions to be defined:
\begin{itemize}
\item $\mathsf{freshId} : \mathcal{N} \to \mathsf{Id}$ is the function that
generates a fresh identifier for a node.

\item $\mathsf{freshName} : \mathcal{N} \to \mathsf{Name}$ is the function that
generates a fresh name for an engine.

\item $\execEffect$ is the function that executes effects as defined in
\Cref{def:effects} and outputs the new system state after the effect has been
executed.
\end{itemize}

We now extend the functions defined in \Cref{rem:node-functions}. The
node-specific functions $\alpha_i$ and $\mathsf{mailboxOf}_i$ can be used to
construct global versions $\alpha$ and $\mathsf{mailboxOf}$ that operate over
all engines $\mathcal{E}$ in $\kappa$, functions of type $\mathcal{E} \to
\mathsf{Addr}$. These global functions are defined by combining the
node-specific functions across all nodes, allowing us to look up addresses for
any engine in the system regardless of which node it resides on.

\section{Operational Semantics}\label{sec:semantics}

We formalise the reduction rules for our operational semantics through the relation $\kappa_n \to \kappa_{n+1}$, where $\kappa_n \equiv (\alpha_n, \mu_n)$ represents the system state. Mailbox engines are modeled as sets for notational convenience, with message operations handled through message passing.

\subsection{System Transition Rules}

\begin{definition}[\textsf{Send}]\label{rule:send}
Processing engine $e$ sends message $v$ to target $e_{\mathsf{target}}$:
\begin{equation}\label{eq:send}
\frac{
  \alpha(e) \equiv \langle b, m_e, r, \mathsf{cfg}, \mathsf{env} \rangle \quad
  m_{\mathsf{target}} \equiv \mathsf{mailboxOf}(e_{\mathsf{target}}) \quad
  v : \mathsf{Msg}
}{
  (\alpha, \mu) \to (\alpha, \mu \cup \{ \langle e, m_{\mathsf{target}}, v \rangle \})
}
\end{equation}
\end{definition}

\begin{definition}[\textsf{Append}]\label{rule:append}
Mailbox engine $e$ enqueues message $v$:
\begin{equation}\label{eq:append}
\frac{
  \alpha(e) \equiv q \quad
  \langle s, e, v \rangle \in \mu
}{
  (\alpha, \mu) \to (\alpha[e \mapsto q \cup \{v\}], \mu \setminus \{\langle s, e, v \rangle\})
}
\end{equation}
\end{definition}

\begin{definition}[\textsf{Pull}]\label{rule:pull}
Processing engine retrieves message from mailbox:
\begin{equation}\label{eq:pull}
\frac{
  \alpha(e) \equiv \langle b, m_e, \mathsf{ready}, \mathsf{cfg}, \mathsf{env} \rangle \quad
  \alpha(m_e) \equiv \{v\} \cup q
}{
  (\alpha, \mu) \to (\alpha', \mu)
}
\end{equation}
where $\alpha' \equiv \alpha[e \mapsto \langle b, m_e, \mathsf{busy}, \mathsf{cfg}, \mathsf{env} \rangle] \cup \{ m_e \mapsto q \}$.
\end{definition}

\begin{definition}[\textsf{Clean}]\label{rule:clean}
Remove terminated engine:
\begin{equation}\label{eq:clean}
\frac{
  \alpha(e) \equiv \langle b, m_e, \mathsf{terminated}, \mathsf{cfg}, \mathsf{env} \rangle
}{
  (\alpha, \mu) \to (\alpha \setminus \{e\}, \mu)
}
\end{equation}
\end{definition}

\begin{definition}[\textsf{Process}]\label{rule:process}
Engine processes message:
\begin{equation}\label{eq:process}
\frac{
  \alpha(e) \equiv \langle b, m_e, \mathsf{busy}, \mathsf{cfg}, \mathsf{env} \rangle \quad
  b, \mathsf{env} \vdash v \Downarrow E \quad
  \execEffect(E, e, (\alpha, \mu)) \equiv (\alpha', \mu', r')
}{
  (\alpha, \mu) \to (\alpha', \mu')
}
\end{equation}
where $r' \equiv \mathsf{ready}$ if $r' \neq \mathsf{terminated}$.
\end{definition}

\subsection{Effect Execution Rules}

The function $\execEffect(E, a, (\alpha, \mu))$ is defined recursively:

\begin{definition}[\textsf{E-sendMsg}]\label{rule:e-send}
\begin{equation}\label{eq:e-send}
\frac{
  \mathsf{msg} \equiv \langle e_{\mathsf{sender}}, e_{\mathsf{target}}, v \rangle \quad
  m_{\mathsf{target}} \equiv \mathsf{mailboxOf}(e_{\mathsf{target}})
}{
  \execEffect(\mathsf{sendMsg}(\mathsf{msg}), a, (\alpha, \mu)) \equiv ((\alpha, \mu \cup \{ \mathsf{msg} \}), \mathsf{ready})
}
\end{equation}
\end{definition}

\begin{definition}[\textsf{E-updateEnv}]\label{rule:e-update}
\begin{equation}\label{eq:e-update}
\frac{
  \alpha(e) \equiv \langle b, m_e, r, \mathsf{cfg}, \mathsf{env} \rangle \quad
  \mathsf{env}^{*} : \mathsf{ofType}(\mathsf{env})
}{
  \execEffect(\mathsf{updateEnv}(\mathsf{env}^{*}), e, (\alpha, \mu)) \equiv (\alpha[e \mapsto \langle b, m_e, r, \mathsf{cfg}, \mathsf{env}^{*} \rangle], \mathsf{ready})
}
\end{equation}
\end{definition}

\begin{definition}[\textsf{E-spawnEngine}]\label{rule:e-spawn}
Create new engine:
\begin{equation}\label{eq:e-spawn}
\frac{
  n \equiv \mathsf{freshId}(\alpha, \mu) \quad
  m \equiv \mathsf{freshId}(\alpha, \mu) \quad
  \alpha(e) \equiv \langle b, m_e, r, \mathsf{cfg}, \mathsf{env} \rangle
}{
  \execEffect(\mathsf{spawnEngine}(e), a, (\alpha, \mu)) \equiv (\alpha \cup \{ n \mapsto \langle b, m, \mathsf{ready}, \mathsf{cfg}, \mathsf{env} \rangle \}, \mathsf{ready})
}
\end{equation}
\end{definition}

\begin{definition}[\textsf{E-chain}]\label{rule:e-chain}
Sequential composition:

\begin{equation}\label{eq:e-chain}
\frac{
  \execEffect(e_1, a, (\alpha, \mu)) \equiv (\alpha_1, \mu_1, s_1) \quad
  \execEffect(\mathsf{chain}(e_2; \dots; e_n), a, (\alpha_1, \mu_1)) \equiv (\alpha', \mu', s')
}{
  \execEffect(\mathsf{chain}(e_1; \dots; e_n), e, (\alpha, \mu)) \equiv (\alpha', \mu', s')
}
\end{equation}
\end{definition}

The complete operational semantics comprises three components:
\begin{enumerate}
\item Configuration transitions via Rules \textsf{Send}--\textsf{Process}
\item Effect execution through Rules \textsf{E-sendMsg}--\textsf{E-chain}
\item Type safety through Definition~\ref{def:effects}
\end{enumerate}

This system formalises message passing (via Rules~\ref{rule:send}--\ref{rule:append}), message processing (Rule~\ref{rule:process}), and effect-driven state changes as defined in Section~\ref{sec:semantics}.

\section{Related Work}\label{sec:related}

Prior work has explored decoupling mailboxes from actors, primarily through
techniques like work-queues where actors pull messages when ready to process
them. Libraries like WeChaty in JavaScript have implemented similar mailbox
separation patterns. However, our work is novel in treating mailboxes as
first-class primitives within the actor model itself.

% \section{Concluding remarks}\label{sec:conclusion}

\section{Acknowledgements}
The inspiration for this work comes from a discussion initiated by Jeremy Ornelas
and Christophe De Troyer, and continued by Tobias Heindel. Thank them for their
feedback and their questions that motivated this work.

\bibliography{art.bib,ref.bib}

\appendix

% \section{Some extra stuff}
% The lean formalization of this work is available at
% \url{https://github.com/anoma/ART-mailboxes-actors}.


\end{document}
